\chapter{Semantyka Kripkego i pełność logik modalnych}
\label{ch:ncl-modal-kripke}

W poprzednim rozdziale wprowadziliśmy normalne
logiki modalne, ich semantykę algebraiczną
oraz przykłady interpretacji relacyjnej.
Teraz badamy systematycznie zależność między
budową relacji dostępności i prawdziwością formuł.
Rozróżniamy własności pojedynczego modelu,
ważność na ramie i ważność w klasie ram.
W szczególności nie zakładamy, że każda normalna
logika modalna jest pełna względem pewnej klasy ram.

\section{Ramy, modele i relacja spełniania}

\begin{definicja}[Rama i model Kripkego]
\label{def:ncl-kripke-frame-model}
Ramą Kripkego nazywamy parę $F=(W,R)$,
gdzie $W\ne\varnothing$ jest zbiorem stanów
(światów), a $R\subseteq W\times W$ jest
relacją dostępności. Piszemy
\[
 R(w)=\{z\in W:wRz\}.
\]
Model Kripkego to trójka $M=(W,R,V)$,
gdzie $(W,R)$ jest ramą, a
$V:\mathrm{Var}\to\mathcal P(W)$ jest
wartościowaniem zmiennych zdaniowych.
Model z wyróżnionym stanem to para $(M,w)$,
gdzie $w\in W$.
\end{definicja}

Nie wymagamy zwrotności, przechodniości
ani symetryczności relacji $R$.
Strzałka $w\to z$ na diagramie oznacza
$wRz$, nie zaś implikację logiczną.
W odróżnieniu od intuicjonistycznych modeli
Kripkego z rozdziału~\ref{ch:ncl-intuitionistic},
wartościowanie nie musi być monotoniczne
względem relacji dostępności.

\begin{definicja}[Spełnianie formuły]
\label{def:ncl-kripke-satisfaction}
Dla języka z Definicji~\ref{def:ncl-modal-language}
definiujemy rekurencyjnie relację $M,w\Vdash A$:
\[
\begin{aligned}
 M,w\Vdash p &\iff w\in V(p),\\
 M,w&\not\Vdash\bot,\\
 M,w\Vdash A\wedge B
   &\iff M,w\Vdash A\ \text{oraz}\ M,w\Vdash B,\\
 M,w\Vdash A\vee B
   &\iff M,w\Vdash A\ \text{lub}\ M,w\Vdash B,\\
 M,w\Vdash A\to B
   &\iff M,w\not\Vdash A\ \text{lub}\ M,w\Vdash B,\\
 M,w\Vdash\Box A
   &\iff \forall z\in R(w)\quad M,z\Vdash A.
\end{aligned}
\]
Z definicji skrótów wynikają klauzule
\[
\begin{aligned}
 M,w\Vdash\neg A&\iff M,w\not\Vdash A,\\
 M,w\Vdash\Diamond A&\iff
 \exists z\in R(w)\quad M,z\Vdash A.
\end{aligned}
\]
\end{definicja}
Pierwsza z ostatnich klauzul wynika z faktu,
że $\bot$ nie jest spełnione w żadnym stanie.
Druga wynika z definicji $\Diamond A=\neg\Box\neg A$
oraz klasycznego zaprzeczenia kwantyfikatora ogólnego.
W szczególności $\Box A$ jest prawdziwa w stanie
bez następców niezależnie od $A$, natomiast
$\Diamond A$ jest w takim stanie fałszywa.
Wartości formuł są dwuwartościowe, choć
prawdziwość może zależeć od stanu.

\begin{przyklad}[Prawdziwość a dostępność]
\label{ex:ncl-kripke-three-states}
Niech $W=\{r,s,t\}$, $R=\{(r,s),(s,t)\}$
oraz $V(p)=\{s\}$. Innych krawędzi nie ma.
\begin{center}
\begin{tikzpicture}[>=Stealth,node distance=22mm]
\node[draw,circle] (r) {$r$};
\node[draw,circle,fill=blue!15,right=of r] (s) {$s$};
\node[draw,circle,right=of s] (t) {$t$};
\draw[->] (r)--(s);
\draw[->] (s)--(t);
\node[below=4mm of r] {$\neg p$};
\node[below=4mm of s] {$p$};
\node[below=4mm of t] {$\neg p$};
\end{tikzpicture}
\end{center}
W stanie $r$ prawdziwe są $\Box p$ i $\Diamond p$,
choć samo $p$ jest fałszywe.
W stanie $s$ prawdziwe jest $p$, ale $\Box p$
i $\Diamond p$ są fałszywe, ponieważ jedyny
następnik $t$ nie spełnia $p$.
W stanie $t$ zachodzi $\Box p$, lecz nie $\Diamond p$.
W stanie $r$ formuła $\Box\Box p$ jest fałszywa,
gdyż $s$ nie spełnia $\Box p$.
Model rozróżnia zatem prawdziwość punktową,
konieczność i możliwość oraz ilustruje
zawodność schematów T i 4.
\end{przyklad}

\begin{stwierdzenie}[Zbiory prawdziwości]
\label{prop:ncl-kripke-truth-sets}
Dla modelu $M=(W,R,V)$ oznaczmy
\[
 [\![A]\!]_M=\{w\in W:M,w\Vdash A\}.
\]
Przypisanie $A\mapsto[\![A]\!]_M$ jest
wartościowaniem w algebrze kompleksowej relacji
z Przykładu~\ref{ex:ncl-modal-complex-algebra}.
W szczególności
\[
 [\![\Box A]\!]_M=\Box_R[\![A]\!]_M,\qquad
 [\![\Diamond A]\!]_M=\Diamond_R[\![A]\!]_M.
\]
\end{stwierdzenie}
\begin{proof}
Dla zmiennych zbiory prawdziwości są równe
$V(p)$, dla $\bot$ otrzymujemy $\varnothing$.
Klauzule klasyczne dają odpowiednio przecięcie,
sumę, dopełnienie i implikację boolowską.
Z klauzuli konieczności mamy
\[
 w\in[\![\Box A]\!]_M
 \iff R(w)\subseteq[\![A]\!]_M.
\]
To jest definicja $\Box_R$.
Klauzula możliwości daje analogicznie
warunek niepustego przecięcia.
Indukcja po budowie formuły pokazuje,
że otrzymane przypisanie jest dokładnie
rozszerzeniem $V$ do wartościowania algebraicznego.
\end{proof}

\begin{lemat}[Lemat o podstawieniu]
\label{lem:ncl-kripke-substitution}
Niech $\sigma$ będzie podstawieniem modalnym.
Dla modelu $M=(W,R,V)$ określmy model
$M_\sigma=(W,R,V_\sigma)$ przez
$V_\sigma(p)=[\![\sigma(p)]\!]_M$.
Wtedy dla każdej formuły $A$ i każdego $w\in W$
\[
 M,w\Vdash\sigma(A)\iff M_\sigma,w\Vdash A.
\]
\end{lemat}
\begin{proof}
Indukujemy po budowie $A$, jednocześnie
dla wszystkich stanów.
Dla zmiennej jest to definicja $V_\sigma$,
a dla $\bot$ obie strony są fałszywe.
Dla spójników klasycznych stosujemy hipotezę
indukcyjną do argumentów w tym samym stanie.
Dla $A=\Box C$ mamy
\[
\begin{aligned}
 M,w\Vdash\sigma(\Box C)
 &\iff \forall z\in R(w)\quad M,z\Vdash\sigma(C)\\
 &\iff \forall z\in R(w)\quad M_\sigma,z\Vdash C\\
 &\iff M_\sigma,w\Vdash\Box C.
\end{aligned}
\]
To kończy indukcję. Skróty są objęte dowodem
po ich rozwinięciu.
\end{proof}

\section{Ważność i konsekwencja relacyjna}

\begin{definicja}[Ważność i spełnialność]
\label{def:ncl-kripke-validity}
Piszemy $M\Vdash A$, jeśli $M,w\Vdash A$
dla wszystkich $w\in W$.
Piszemy $F\Vdash A$, jeśli $M\Vdash A$
dla każdego modelu $M$ na ramie $F$.
Dla klasy ram $\mathcal C$ piszemy
$\mathcal C\Vdash A$, jeśli $F\Vdash A$
dla wszystkich $F\in\mathcal C$.
Formuła $A$ jest spełnialna w klasie
$\mathcal C$, jeśli istnieje $F\in\mathcal C$,
model $M$ na $F$ oraz stan $w$, w którym
$M,w\Vdash A$.
\end{definicja}
Ważność na ramie wymaga kwantyfikacji po
wszystkich wartościowaniach, nie tylko po
stanach jednego modelu. Z klasycznej negacji
wynika, że $A$ jest ważna w $\mathcal C$
dokładnie wtedy, gdy $\neg A$ nie jest
spełnialna w $\mathcal C$.

\begin{definicja}[Konsekwencja lokalna i globalna]
\label{def:ncl-kripke-consequence}
Niech $\mathcal C$ będzie klasą ram.
Dla zbioru formuł $\Gamma$ i formuły $A$
określamy dwa pojęcia konsekwencji:
\[
 \Gamma\models_{\mathcal C}^{\mathrm{lok}} A
\]
oznacza, że dla każdego modelu $M$ na ramie
z $\mathcal C$ i każdego stanu $w$,
jeżeli $M,w\Vdash C$ dla każdego $C\in\Gamma$,
to $M,w\Vdash A$.
Zapis
\[
 \Gamma\models_{\mathcal C}^{\mathrm{glob}} A
\]
oznacza, że dla każdego takiego modelu $M$,
jeżeli $M\Vdash C$ dla każdego $C\in\Gamma$,
to $M\Vdash A$.
\end{definicja}
W pierwszym przypadku wszystkie przesłanki
muszą zachodzić w jednym stanie, w drugim
w każdym stanie ustalonego modelu.
W obu przypadkach wartościowanie jest wspólne
dla wszystkich przesłanek i wniosku.

\begin{stwierdzenie}[Podstawowe własności konsekwencji]
\label{prop:ncl-kripke-consequence-laws}
Oba rodzaje konsekwencji są zwrotne,
monotoniczne, spełniają regułę cięcia
i są strukturalne. Ponadto
\[
 \Gamma\models_{\mathcal C}^{\mathrm{lok}}A
 \Longrightarrow
 \Gamma\models_{\mathcal C}^{\mathrm{glob}}A,
 \qquad
 \varnothing\models_{\mathcal C}^{\mathrm{lok}}A
 \iff \varnothing\models_{\mathcal C}^{\mathrm{glob}}A
 \iff\mathcal C\Vdash A.
\]
Dla skończonego $\Gamma$ zachodzi
\[
 \Gamma\models_{\mathcal C}^{\mathrm{lok}}A
 \iff\mathcal C\Vdash
        \left(\bigwedge_{C\in\Gamma}C\right)\to A.
\]
Pusta koniunkcja oznacza $\top$.
\end{stwierdzenie}
\begin{proof}
Zwrotność i monotoniczność wynikają
bezpośrednio z definicji. Dla cięcia,
jeśli wszystkie przesłanki z $\Gamma$ są
spełnione, to każda konsekwencja z $\Delta$
jest spełniona, a zatem również konsekwencja
$A$ zbioru $\Delta$. W przypadku lokalnym
wszystkie trzy kroki odnoszą się do tego
samego stanu, w globalnym do całego modelu.
Strukturalność wynika z
Lematu~\ref{lem:ncl-kripke-substitution}:
model $M_\sigma$ ma tę samą ramę co $M$,
więc należy do tej samej rozważanej klasy modeli.
Jeżeli przesłanki są prawdziwe globalnie,
to są prawdziwe lokalnie w każdym stanie.
Zastosowanie konsekwencji lokalnej w każdym
stanie daje konsekwencję globalną.
Dla pustego zbioru przesłanek oba warunki
sprowadzają się do ważności w klasie.
Ostatnia równoważność wynika z klasycznych
klauzul koniunkcji i implikacji w każdym stanie.
\end{proof}
Nie wpisujemy na tej liście skończoności
konsekwencji: dla dowolnej klasy ram własność
ta nie musi zachodzić. Będziemy ją analizować
wraz z twierdzeniami o silnej pełności.

\begin{przyklad}[Konsekwencja globalna nie musi być lokalna]
\label{ex:ncl-kripke-global-not-local}
Dla każdej klasy ram $\mathcal C$ mamy
$p\models_{\mathcal C}^{\mathrm{glob}}\Box p$:
jeśli $p$ jest prawdziwe w każdym stanie,
to jest prawdziwe w każdym następniku
dowolnego stanu.
Niech $\mathcal C$ zawiera ramę
dwuelementową z relacją pełną.
Dla $V(p)=\{r\}$ w stanie $r$ prawdziwe jest
$p$, lecz $\Box p$ jest fałszywa,
ponieważ drugi stan $s$ nie spełnia $p$.
Stąd $p\not\models_{\mathcal C}^{\mathrm{lok}}\Box p$.
Kontrprzykład działa nawet w klasie ram S5.
\end{przyklad}

\begin{twierdzenie}[Poprawność relacyjna normalnych logik]
\label{thm:ncl-kripke-soundness}
Niech $L$ będzie normalną logiką modalną
i niech $\mathcal C$ będzie klasą ram,
na których ważne są wszystkie formuły $L$.
Wtedy
\[
\begin{aligned}
 \Gamma\vdash_L^{\mathrm{lok}} A
 &\Longrightarrow
 \Gamma\models_{\mathcal C}^{\mathrm{lok}}A,\\
 \Gamma\vdash_L^{\mathrm{glob}} A
 &\Longrightarrow
 \Gamma\models_{\mathcal C}^{\mathrm{glob}}A.
\end{aligned}
\]
System K jest poprawny względem klasy
wszystkich niepustych ram Kripkego.
\end{twierdzenie}
\begin{proof}
Dla pierwszych dwóch tez stosujemy indukcję
po dowodzie z
Definicji~\ref{def:ncl-modal-local-global}.
Każda przesłanka jest prawdziwa z założenia,
każde twierdzenie $L$ z ważności na ramie.
Modus ponens zachowuje prawdziwość w stanie,
a więc również prawdziwość globalną.
W dowodzie globalnym reguła konieczności
zachowuje prawdziwość globalną: jeżeli
$M\Vdash C$, to w każdym stanie $w$
wszyscy jego następnicy spełniają $C$,
czyli $M\Vdash\Box C$.
W dowodzie lokalnym tej dodatkowej
reguły dla przesłanek nie ma.

Dla ostatniej tezy sprawdzamy aksjomaty K.
Klasyczne tautologie są prawdziwe w każdym
stanie ze względu na klasyczne klauzule
spójników. Jeżeli w stanie $w$ zachodzą
$\Box(A\to B)$ oraz $\Box A$, to każdy
$z\in R(w)$ spełnia zarówno $A\to B$, jak i $A$,
a więc także $B$. Zatem $w$ spełnia $\Box B$,
co daje ważność aksjomatu K.
Obie reguły dowodzenia twierdzeń zachowują
ważność w każdym modelu na każdej ramie,
jak wykazano wyżej. Indukcja po dowodach
K daje ważność wszystkich twierdzeń K,
co pozwala zastosować pierwsze dwie tezy.
\end{proof}

\begin{wniosek}[Logika klasy ram]
\label{cor:ncl-kripke-frame-logic}
Dla dowolnej klasy ram $\mathcal C$ zbiór
\[
 \operatorname{Log}(\mathcal C)
 =\{A:\mathcal C\Vdash A\}
\]
jest normalną logiką modalną.
Jeżeli $\mathcal C$ jest niepusta,
to logika ta jest niesprzeczna.
\end{wniosek}
\begin{proof}
Poprawność K daje zawieranie aksjomatów K.
Modus ponens i konieczność zachowują
ważność, a podstawienie zachowuje ją
z Lematu~\ref{lem:ncl-kripke-substitution}.
Jeśli istnieje rama z $\mathcal C$,
to ma ona co najmniej jeden stan.
Formuła $\bot$ jest w nim fałszywa,
więc nie należy do $\operatorname{Log}(\mathcal C)$.
Dla pustej klasy wszystkie formuły są
ważne próżniowo, co daje logikę sprzeczną.
\end{proof}
